\begin{abstract}
Autoregressive large language model (LLM) serving systems that schedule requests
first-come-first-served suffer from head-of-line (HOL) blocking: because the
output length of a request is unknown until decoding finishes, a single long
generation can delay many short ones behind it in the same batch. Recent
learning-to-rank (LTR) schedulers address this by training a small predictor to
estimate a request's relative output length and using that score to reorder
the admission queue, but the learned score alone can be miscalibrated,
particularly for prompts that fall outside the predictor's training
distribution. This work extends an existing OPT-125M LTR predictor with a
lightweight, non-learned secondary ranking signal, the tokenized prompt length,
combined with the predictor's score through a weighted composite formula. The
secondary signal requires no additional model inference and is available for
every request at negligible cost.

We integrate the composite scheduler into vLLM and evaluate it against
first-come-first-served, a classification-based baseline, and the original
LTR scheduler on two 1{,}000-request subsets of LMSYS-Chat-1M served with
Llama-3.1-8B-Instruct on a single NVIDIA A100 GPU. The composite scheduler
achieves the highest mean throughput of the four policies on both the
in-distribution and out-of-distribution test sets, improving throughput by
roughly 31.7\% over the original LTR scheduler on the in-distribution set, and
modestly improves ranking quality under distribution shift, raising Kendall's
tau from 0.3445 to 0.3457 on the out-of-distribution set.
\end{abstract}
